\section{Advanced Control Flow and Scope Management}
\label{sec:control-flow-and-scope-management}

% ─────────────────────────────────────────────
\subsection{The Parent Operator \texttt{\^{}} and Scope Hierarchy}
\label{sec:parent-operator}

\subsubsection{Scope in QLang}
In QLang, every block of code creates a new scope level, forming a linked chain of
scopes. Reading a variable walks this chain upward automatically. Writing to a
variable in an outer scope requires the parent operator \texttt{\^{}}.

\subsubsection{How \texttt{\^{}} Works}
The parent operator \texttt{\^{}} explicitly selects a scope one level above the current
one. Multiple carets chain upward, one level per caret:

\begin{itemize}
    \item \texttt{\^{}x} accesses \texttt{x} in the parent scope
    \item \texttt{\^{}\^{}x} accesses \texttt{x} in the grandparent scope
    \item \texttt{\^{}\^{}\^{}x} accesses \texttt{x} three levels up
\end{itemize}

Once the target scope is reached, the variable is looked up only in that
specific scope; the automatic upward walk does not apply at that point.
The operator supports all standard operations: assignment, compound assignment,
increment and decrement, as well as indexed access on arrays.

The following example demonstrates reading and writing through the scope hierarchy:
% Auto-generated by docs/tools/highlight.py - do not edit
\begin{codebox}
(*@\textcolor{codebox_type}{\textbf{num}}@*) (*@\textcolor{black}{x}@*) (*@\textcolor{codebox_operator}{=}@*) (*@\textcolor{codebox_number}{1}@*)(*@\textcolor{codebox_operator}{;}@*)
(*@\textcolor{codebox_control}{\textbf{if}}@*) (*@\textcolor{codebox_operator}{(}@*)(*@\textcolor{codebox_number}{1}@*)(*@\textcolor{codebox_operator}{)}@*) (*@\textcolor{codebox_operator}{\{}@*)
    (*@\textcolor{codebox_type}{\textbf{num}}@*) (*@\textcolor{black}{x}@*) (*@\textcolor{codebox_operator}{=}@*) (*@\textcolor{codebox_number}{2}@*)(*@\textcolor{codebox_operator}{;}@*)
    (*@\textcolor{codebox_control}{\textbf{if}}@*) (*@\textcolor{codebox_operator}{(}@*)(*@\textcolor{codebox_number}{1}@*)(*@\textcolor{codebox_operator}{)}@*) (*@\textcolor{codebox_operator}{\{}@*)
        (*@\textcolor{codebox_string}{println}@*)(*@\textcolor{codebox_operator}{(}@*)(*@\textcolor{codebox_string}{\string"\%d\string"}@*) (*@\textcolor{codebox_operator}{\%}@*) (*@\textcolor{codebox_operator}{\^{ }}@*)(*@\textcolor{black}{x}@*)(*@\textcolor{codebox_operator}{)}@*)(*@\textcolor{codebox_operator}{;}@*)   (*@\textcolor{codebox_comment}{// 2  - reads parent x}@*)
        (*@\textcolor{codebox_operator}{\^{ }}@*)(*@\textcolor{black}{x}@*) (*@\textcolor{codebox_operator}{=}@*) (*@\textcolor{codebox_number}{7}@*)(*@\textcolor{codebox_operator}{;}@*)
        (*@\textcolor{codebox_string}{println}@*)(*@\textcolor{codebox_operator}{(}@*)(*@\textcolor{codebox_string}{\string"\%d\string"}@*) (*@\textcolor{codebox_operator}{\%}@*) (*@\textcolor{codebox_operator}{\^{ }}@*)(*@\textcolor{black}{x}@*)(*@\textcolor{codebox_operator}{)}@*)(*@\textcolor{codebox_operator}{;}@*)   (*@\textcolor{codebox_comment}{// 7  - assigns to parent x}@*)
        (*@\textcolor{codebox_operator}{(}@*)(*@\textcolor{codebox_operator}{\^{ }}@*)(*@\textcolor{black}{x}@*)(*@\textcolor{codebox_operator}{)}@*)++(*@\textcolor{codebox_operator}{;}@*)
        (*@\textcolor{codebox_string}{println}@*)(*@\textcolor{codebox_operator}{(}@*)(*@\textcolor{codebox_string}{\string"\%d\string"}@*) (*@\textcolor{codebox_operator}{\%}@*) (*@\textcolor{codebox_operator}{\^{ }}@*)(*@\textcolor{black}{x}@*)(*@\textcolor{codebox_operator}{)}@*)(*@\textcolor{codebox_operator}{;}@*)   (*@\textcolor{codebox_comment}{// 8}@*)
        (*@\textcolor{codebox_operator}{\^{ }}@*)(*@\textcolor{black}{x}@*) += (*@\textcolor{codebox_number}{3}@*)(*@\textcolor{codebox_operator}{;}@*)
        (*@\textcolor{codebox_string}{println}@*)(*@\textcolor{codebox_operator}{(}@*)(*@\textcolor{codebox_string}{\string"\%d\string"}@*) (*@\textcolor{codebox_operator}{\%}@*) (*@\textcolor{codebox_operator}{\^{ }}@*)(*@\textcolor{black}{x}@*)(*@\textcolor{codebox_operator}{)}@*)(*@\textcolor{codebox_operator}{;}@*)   (*@\textcolor{codebox_comment}{// 11}@*)
        (*@\textcolor{codebox_string}{println}@*)(*@\textcolor{codebox_operator}{(}@*)(*@\textcolor{codebox_string}{\string"\%d\string"}@*) (*@\textcolor{codebox_operator}{\%}@*) (*@\textcolor{codebox_operator}{\^{ }}@*)(*@\textcolor{codebox_operator}{\^{ }}@*)(*@\textcolor{black}{x}@*)(*@\textcolor{codebox_operator}{)}@*)(*@\textcolor{codebox_operator}{;}@*)  (*@\textcolor{codebox_comment}{// 1  - reads grandparent x}@*)
    (*@\textcolor{codebox_operator}{\}}@*)
(*@\textcolor{codebox_operator}{\}}@*)
\end{codebox}


Array indexing works the same way. The parent operator applies to the variable,
and the index is evaluated afterward:
% Auto-generated by docs/tools/highlight.py - do not edit
\begin{codebox}
(*@\textcolor{codebox_type}{\textbf{num}}@*) (*@\textcolor{black}{arr}@*)(*@\textcolor{codebox_operator}{[}@*)(*@\textcolor{codebox_number}{2}@*)(*@\textcolor{codebox_operator}{]}@*) (*@\textcolor{codebox_operator}{=}@*) (*@\textcolor{codebox_operator}{[}@*)(*@\textcolor{codebox_number}{10}@*)(*@\textcolor{codebox_operator}{,}@*) (*@\textcolor{codebox_number}{20}@*)(*@\textcolor{codebox_operator}{]}@*)(*@\textcolor{codebox_operator}{;}@*)
(*@\textcolor{codebox_control}{\textbf{if}}@*) (*@\textcolor{codebox_operator}{(}@*)(*@\textcolor{codebox_number}{1}@*)(*@\textcolor{codebox_operator}{)}@*) (*@\textcolor{codebox_operator}{\{}@*)
    (*@\textcolor{codebox_control}{\textbf{if}}@*) (*@\textcolor{codebox_operator}{(}@*)(*@\textcolor{codebox_number}{1}@*)(*@\textcolor{codebox_operator}{)}@*) (*@\textcolor{codebox_operator}{\{}@*)
        (*@\textcolor{codebox_string}{println}@*)(*@\textcolor{codebox_operator}{(}@*)(*@\textcolor{codebox_string}{\string"\%d\string"}@*) (*@\textcolor{codebox_operator}{\%}@*) (*@\textcolor{codebox_operator}{\^{ }}@*)(*@\textcolor{codebox_operator}{\^{ }}@*)(*@\textcolor{black}{arr}@*)(*@\textcolor{codebox_operator}{[}@*)(*@\textcolor{codebox_number}{0}@*)(*@\textcolor{codebox_operator}{]}@*)(*@\textcolor{codebox_operator}{)}@*)(*@\textcolor{codebox_operator}{;}@*)  (*@\textcolor{codebox_comment}{// 10}@*)
        (*@\textcolor{codebox_string}{println}@*)(*@\textcolor{codebox_operator}{(}@*)(*@\textcolor{codebox_string}{\string"\%d\string"}@*) (*@\textcolor{codebox_operator}{\%}@*) (*@\textcolor{codebox_operator}{\^{ }}@*)(*@\textcolor{codebox_operator}{\^{ }}@*)(*@\textcolor{black}{arr}@*)(*@\textcolor{codebox_operator}{[}@*)(*@\textcolor{codebox_number}{1}@*)(*@\textcolor{codebox_operator}{]}@*)(*@\textcolor{codebox_operator}{)}@*)(*@\textcolor{codebox_operator}{;}@*)  (*@\textcolor{codebox_comment}{// 20}@*)
    (*@\textcolor{codebox_operator}{\}}@*)
(*@\textcolor{codebox_operator}{\}}@*)
\end{codebox}


% ─────────────────────────────────────────────
\subsection{Advanced Array Structures and Indexing}
\label{sec:array-indexing}

QLang supports three indexing modes for arrays: simple indexing, range indexing,
and list indexing. All three resolve negative indices automatically against the
container length.

\subsubsection{Simple Indexing and Negative Indices}
A single integer index selects one element. Negative values count from the end
of the array: an index \texttt{i} where \texttt{i < 0} is resolved as
\texttt{i + length} before access.

% Auto-generated by docs/tools/highlight.py - do not edit
\begin{codebox}
(*@\textcolor{codebox_type}{\textbf{num}}@*) (*@\textcolor{black}{arr}@*)(*@\textcolor{codebox_operator}{[}@*)(*@\textcolor{codebox_number}{5}@*)(*@\textcolor{codebox_operator}{]}@*) (*@\textcolor{codebox_operator}{=}@*) (*@\textcolor{codebox_operator}{[}@*)(*@\textcolor{codebox_number}{10}@*)(*@\textcolor{codebox_operator}{,}@*) (*@\textcolor{codebox_number}{20}@*)(*@\textcolor{codebox_operator}{,}@*) (*@\textcolor{codebox_number}{30}@*)(*@\textcolor{codebox_operator}{,}@*) (*@\textcolor{codebox_number}{40}@*)(*@\textcolor{codebox_operator}{,}@*) (*@\textcolor{codebox_number}{50}@*)(*@\textcolor{codebox_operator}{]}@*)(*@\textcolor{codebox_operator}{;}@*)

(*@\textcolor{codebox_string}{println}@*)(*@\textcolor{codebox_operator}{(}@*)(*@\textcolor{black}{arr}@*)(*@\textcolor{codebox_operator}{[}@*)(*@\textcolor{codebox_number}{0}@*)(*@\textcolor{codebox_operator}{]}@*)(*@\textcolor{codebox_operator}{)}@*)(*@\textcolor{codebox_operator}{;}@*)   (*@\textcolor{codebox_comment}{// 10}@*)

(*@\textcolor{codebox_string}{println}@*)(*@\textcolor{codebox_operator}{(}@*)(*@\textcolor{black}{arr}@*)(*@\textcolor{codebox_operator}{[}@*)(*@\textcolor{codebox_operator}{-}@*)(*@\textcolor{codebox_number}{1}@*)(*@\textcolor{codebox_operator}{]}@*)(*@\textcolor{codebox_operator}{)}@*)(*@\textcolor{codebox_operator}{;}@*)  (*@\textcolor{codebox_comment}{// 50}@*)

(*@\textcolor{codebox_string}{println}@*)(*@\textcolor{codebox_operator}{(}@*)(*@\textcolor{black}{arr}@*)(*@\textcolor{codebox_operator}{[}@*)(*@\textcolor{codebox_operator}{-}@*)(*@\textcolor{codebox_number}{2}@*)(*@\textcolor{codebox_operator}{]}@*)(*@\textcolor{codebox_operator}{)}@*)(*@\textcolor{codebox_operator}{;}@*)  (*@\textcolor{codebox_comment}{// 40}@*)
\end{codebox}


\subsubsection{Range Indexing}
A range index \texttt{start..end} selects a contiguous subarray. The range is
left-inclusive and right-exclusive: element at \texttt{start} is included,
element at \texttt{end} is not. Negative values are resolved the same way as
in simple indexing.

% Auto-generated by docs/tools/highlight.py - do not edit
\begin{codebox}
(*@\textcolor{codebox_type}{\textbf{num}}@*) (*@\textcolor{black}{arr}@*)(*@\textcolor{codebox_operator}{[}@*)(*@\textcolor{codebox_number}{5}@*)(*@\textcolor{codebox_operator}{]}@*) (*@\textcolor{codebox_operator}{=}@*) (*@\textcolor{codebox_operator}{[}@*)(*@\textcolor{codebox_number}{10}@*)(*@\textcolor{codebox_operator}{,}@*) (*@\textcolor{codebox_number}{20}@*)(*@\textcolor{codebox_operator}{,}@*) (*@\textcolor{codebox_number}{30}@*)(*@\textcolor{codebox_operator}{,}@*) (*@\textcolor{codebox_number}{40}@*)(*@\textcolor{codebox_operator}{,}@*) (*@\textcolor{codebox_number}{50}@*)(*@\textcolor{codebox_operator}{]}@*)(*@\textcolor{codebox_operator}{;}@*)
(*@\textcolor{codebox_string}{println}@*)(*@\textcolor{codebox_operator}{(}@*)(*@\textcolor{black}{arr}@*)(*@\textcolor{codebox_operator}{[}@*)(*@\textcolor{codebox_number}{1}@*)..(*@\textcolor{codebox_number}{3}@*)(*@\textcolor{codebox_operator}{]}@*)(*@\textcolor{codebox_operator}{)}@*)(*@\textcolor{codebox_operator}{;}@*)   (*@\textcolor{codebox_comment}{// [20, 30]}@*)
(*@\textcolor{codebox_string}{println}@*)(*@\textcolor{codebox_operator}{(}@*)(*@\textcolor{black}{arr}@*)(*@\textcolor{codebox_operator}{[}@*)(*@\textcolor{codebox_number}{0}@*)..(*@\textcolor{codebox_operator}{-}@*)(*@\textcolor{codebox_number}{1}@*)(*@\textcolor{codebox_operator}{]}@*)(*@\textcolor{codebox_operator}{)}@*)(*@\textcolor{codebox_operator}{;}@*)  (*@\textcolor{codebox_comment}{// [10, 20, 30, 40]}@*)
\end{codebox}

% ─────────────────────────────────────────────
\subsubsection{List Indexing}
A list index \texttt{[[a, b, c]]} selects elements at the given positions in
arbitrary order. The double bracket syntax distinguishes list indexing from
simple indexing. This is particularly useful for selecting non-contiguous
elements or reordering array contents.

\input{chapter-4-control-flow-and-scope-management/qlang/fancy-indexing.tex}

An existing array can also be passed as the index list directly:

\input{chapter-4-control-flow-and-scope-management/qlang/index-list-array.tex}
% ─────────────────────────────────────────────
\subsection{...}
